\section{Step 12 Copy the incoming state into working registers}\label{sha-s12}

Load the eight state words into \texttt{s0} through \texttt{s7}, representing a through h. Keep the incoming state in memory for feed-forward at the end of the block. The supplied frame will restore these preserved registers on every return; while the function runs, they are its working storage. A signed word load preserves the low-word pattern required by word operations.

\subsection{Write and inspect}
Replace only the temporary instructions between \texttt{\# BEGIN comp\_load} and \texttt{\# END comp\_load}. Keep the supplied labels outside that region.

\Needspace{12\baselineskip}
\MARSCodeLabel{MIPS assembly}
\begin{minted}[samepage=true]{mars-idea}
    lw $s0, 0($t9)
    lw $s1, 4($t9)
    lw $s2, 8($t9)
    lw $s3, 12($t9)
    lw $s4, 16($t9)
    lw $s5, 20($t9)
    lw $s6, 24($t9)
    lw $s7, 28($t9)
\end{minted}

Use recipe \texttt{s12-1} with the cumulative program through this edit. The exact recipe names the entry, prepared input and checkpoint. Observe the checkpoint before continuing to its temporary or normal return.

\subsection{Expected observation}
At the first-round boundary, \texttt{s0} through \texttt{s7} contain the initial working words and caller state is unchanged.

\Needspace{9\baselineskip}
\subsection{Check your understanding}
Why retain the incoming state words in memory?

\begin{itemize}
\item \textbf{A} They are used again in the final feed-forward sums.
\item \textbf{B} They contain the file path.
\item \textbf{C} They are the schedule destination.
\end{itemize}
